\section{Discussion}
\label{sec:discussion}

\subsection{Key Findings and Their Implications}

\textbf{Trustworthiness rankings are stable latent model properties.} The core empirical result --- high partial $\rho$ for both fairness and adversarial robustness after capability control --- indicates that trustworthiness properties are not test-specific artifacts. For practitioners: model selection decisions based on trustworthiness rankings are more robust than one might have assumed. The relative ordering of models on fairness and robustness benchmarks is preserved when those benchmarks become harder or shift context.

\textbf{The adversarial disruption hypothesis is wrong at the model-ranking level.} At $N=13$ with diverse model families, the disruption disappears. Both adversarial pairs show substantially positive partial $\rho$ with zero rank reversals. Creating harder adversarial evaluation conditions does not reshuffle the model ranking --- the models most robust to adversarial perturbation in easy conditions remain the most robust in hard conditions.

\textbf{Capability control reveals hidden differential structure.} Without MMLU control, raw $\rho$ is near-ceiling for all dimensions and the fairness-robustness comparison is invisible ($\Delta\rho_{\text{raw}} \approx -0.005$). The partial analysis reveals the hidden structure ($\Delta\rho_{\text{partial}} = 0.192$). Evaluators who report only raw cross-benchmark correlations miss dimension-specific information.

\subsection{Limitations}

\textbf{GLUE/AdvGLUE arm structurally unavailable.} GLUE-X \citep{Yang2023GLUEX} evaluates encoder-only PLMs --- zero model overlap with our decoder-only LLMs. We used TrustLLM's OOD robustness task as a proxy. Direct decoder-only LLM evaluation on AdvGLUE would require new data collection.

\textbf{$\Delta\rho = 0.192$ falls below preregistered threshold.} The 0.008 shortfall is within $N=13$ measurement uncertainty. Two contributing factors: three models missing robustness scores reduce power, and $\rho_{\text{ANLI}} = 0.684$ is higher than expected. The directional claim survives; the preregistered confirmation does not. For the Fisher $z$-test validity, $\rho_{\text{fairness}}$ was re-estimated on the $N=13$ subset ($\rho=0.967$) before computing $\Delta\rho=0.192$ and Fisher $z=2.265$ --- ensuring the \citet{Meng1992Comparing} same-sample assumption is met.

\textbf{Causal mechanism unknown.} We proposed adversarial design explains the fairness advantage. This mechanism is falsified. Two alternatives remain viable: (1) \emph{Stable latent bias hypothesis} --- fairness biases more deeply encoded than robustness properties. (2) \emph{Benchmark overlap hypothesis} --- BBQ splits share sufficient item structure that $\rho$ reflects benchmark design, not model mechanism. Item-level Jaccard analysis would distinguish these.

Winogrande sensitivity ($\Delta\rho$ collapses to 0.073 under alternative covariate) suggests robustness $\rho$ values may be partially driven by capability dimensions beyond MMLU. Future work with richer covariates would clarify the $\Delta\rho$ claim.

\textbf{Scope restricted to 2023--2024 model population ($N=16$).} Post-2024 models may exhibit different stability patterns.

\textbf{P3 (instruction-tuning effects) untested.} Whether RLHF alignment improves fairness generalization specifically was deferred. It remains the most practically actionable open question.

\subsection{Broader Impact}

This research provides empirical grounding for an implicit assumption in LLM safety evaluation. The finding that trustworthiness rankings are stable is positive for practitioners relying on benchmark-based model selection. We note that rank stability and absolute performance are distinct --- models maintaining relative ordering on harder benchmarks while all performing worse means adversarial evaluation still reveals genuine capability limits. Stability of rank does not imply adequacy of performance.
